\reading{LTR-17}{Risk Management for Changing Interest Rates: Asset-Liability Management and Duration Techniques}{STRIKE FIRST}{5}
\markboth{LTR-17\ \ ALM and Duration Techniques}{}

\begin{lrkeybox}[Learning objectives — official 2026 spine wording]
\begin{enumerate}[label=\textbf{\alph*.}]
\item Discuss how asset-liability management strategies can help a bank hedge against interest rate risk.
\item Describe interest-sensitive gap management and apply this strategy to maximize a bank's net interest margin.
\item Describe duration gap management and apply this strategy to protect a bank's net worth.
\item Discuss the limitations of interest-sensitive gap management and duration gap management.
\end{enumerate}
{\footnotesize\color{FRMGrey}Source: Peter Rose and Sylvia Hudgins, \textit{Bank Management \& Financial Services}, 9th Edition, Chapter 7.}
\end{lrkeybox}

\begin{lrtrapbox}[Label collision]
The Lecture layer labels this reading \texttt{LO 76.a--d}; the Crash layer uses \texttt{76} for \textbf{LTR-12 Managing Non-Deposit Liabilities}. Renumbered to the spine. At \textbf{SRC 5} this is the second-highest-yield reading in the area, behind only LTR-1.
\end{lrtrapbox}

% =====================================================================
\lo{LTR-17 a}{Discuss how asset-liability management strategies can help a bank hedge against interest rate risk.}

Banks manage risk and profitability through \term{asset-liability management (ALM)}. This considers both a bank's assets (loans, investments) and liabilities (deposits, borrowings), to ensure the interest income earned on assets covers the interest paid on liabilities. Interest rate changes can significantly impact a bank's \textbf{net interest income (NII)} and its \textbf{equity}.

\begin{center}
\small
\begin{tabularx}{\textwidth}{@{}p{4.0cm} X@{}}
\toprule
\rowcolor{FRMSoft}\textbf{ALM strategy} & \textbf{What it does} \\
\midrule
\textbf{Matching maturities} & Aligning the repayment schedules of assets and liabilities. \\
\textbf{Interest rate layering} & Adjusting the mix of fixed- and variable-rate assets and liabilities. \\
\textbf{Off-balance-sheet solutions} & Using instruments such as futures and swaps to hedge against interest rate movements. \\
\bottomrule
\end{tabularx}
\end{center}

\begin{lrkeybox}[The three objectives of asset-liability management]
\begin{enumerate}
\item Management should exercise as much control as possible over the \textbf{volume, mix, and return or cost} of \emph{both} assets and liabilities, in order to achieve the institution's goals.
\item Control over assets must be \textbf{coordinated with} control over liabilities, so that asset management and liability management are internally consistent and \textbf{do not pull against each other}.
\item Revenues and costs arise from \textbf{both sides} of the balance sheet. Management policies need to maximize returns and effectively control costs from supplying services.
\end{enumerate}
\end{lrkeybox}

\begin{center}
\begin{tikzpicture}[
  bx/.style={draw=FRMNavy, fill=PaleNavy, rounded corners=1pt, align=center,
             font=\scriptsize\sffamily, text width=2.5cm, inner sep=4pt},
  bxg/.style={draw=FRMGreen, fill=PaleGreen, rounded corners=1pt, align=center,
             font=\scriptsize\sffamily\bfseries, text width=2.5cm, inner sep=4pt},
  bxo/.style={draw=FRMGold, fill=PaleGold, rounded corners=1pt, align=center,
             font=\scriptsize\sffamily\bfseries, text width=3.0cm, inner sep=4pt},
  ar/.style={-{Stealth[length=4.5pt]}, FRMGold, thick}
]
\node[bx] (src) at (0,0) {Managing the response to changing market interest rates};
\node[bx] (mid) at (4.0,1.0) {Interest revenues\\ Interest costs};
\node[bx] (mid2) at (4.0,-1.0) {Market value of assets\\ Market value of liabilities};
\node[bxg] (nim) at (8.0,1.0) {Net interest margin};
\node[bxg] (nw) at (8.0,-1.0) {Net worth (equity)};
\node[bxo] (out) at (12.2,0) {Investment value, profitability, and risk exposure};
\draw[ar] (src.east) -- (mid.west);
\draw[ar] (src.east) -- (mid2.west);
\draw[ar] (mid) -- (nim);
\draw[ar] (mid2) -- (nw);
\draw[ar] (nim.east) -- (out.north west);
\draw[ar] (nw.east) -- (out.south west);
\end{tikzpicture}
\figcap{Two channels, two tools, and this diagram is the map of the whole reading. Interest \emph{flows} drive the \textbf{net interest margin}, which is what \textbf{interest-sensitive gap management} (objective b) protects. Interest-driven \emph{market values} drive \textbf{net worth}, which is what \textbf{duration gap management} (objective c) protects. Every question in this reading sits on one branch or the other --- identify the branch first.}
\end{center}

\begin{lrkeybox}[Forces determining interest rates]
Whether financial firms are on the supply side or the demand side of the loanable funds market at any given moment, they \textbf{cannot determine the level, or be sure about the trend, of market interest rates}. Most financial managers must be \textbf{price takers, not price makers}, and must accept interest rate levels as given and plan accordingly.

Firms use asset-liability strategies to manage two major kinds of interest rate risk: \textbf{price risk} and \textbf{reinvestment risk}.

Many banks conduct their ALM strategies under the guidance of an \textbf{asset-liability committee (ALCO)}, which meets regularly to manage the firm's interest rate risk and other exposures, centred on the maximization of shareholder wealth, maintaining adequate profitability, and achieving sufficient capitalization.
\end{lrkeybox}

% =====================================================================
\lo{LTR-17 b}{Describe interest-sensitive gap management and apply this strategy to maximize a bank's net interest margin.}

An important goal is to \textbf{insulate profits from the damaging effects of fluctuating interest rates}. To protect profits against adverse rate changes, management seeks to hold the \textbf{net interest margin} fixed.

\begin{lrfmlbox}[Net interest margin]
\[ \text{NIM} \;=\; \frac{\text{NII}}{\text{Earning assets}} \;=\; \frac{\begin{pmatrix}\text{Interest income from}\\ \text{loans and investments}\end{pmatrix} - \begin{pmatrix}\text{Interest expense on deposits}\\ \text{and other borrowed funds}\end{pmatrix}}{\text{Total earning assets}} \]
\vk{NII = net interest income = interest income $-$ interest expense. \textbf{Earning assets} are the bank assets that earn income --- loans and securities. Note the denominator: \emph{earning} assets, not total assets.}
\end{lrfmlbox}

\begin{lrexbox}[Gray Sky Bank's net interest margin]
Gray Sky Bank has total assets of \$50 billion and \textbf{earning assets of \$45 billion}. The bank earned \$4 billion in interest income and paid \$2.2 billion in interest expense. Calculate the bank's NIM.
\tcblower
\textbf{Step 1 --- net interest income.}
\[ \text{NII} = \$4\text{bn} - \$2.2\text{bn} = \$1.8\text{bn} \]

\textbf{Step 2 --- divide by \emph{earning} assets.}
\[ \text{NIM} = \frac{1.8}{45} = \mathbf{4.00\%} \]

\textbf{What this means.} The \$5 billion of non-earning assets --- premises, cash, required reserves --- are deliberately excluded. They consume funding but generate no interest income, so including them would understate the margin the interest-earning book is actually producing.
\end{lrexbox}

\begin{lrtrapbox}[Total assets is the offered denominator]
$1.8 / 50 = 3.60\%$. That is a clean, plausible number produced by one wrong denominator, and the stem supplies the \$50 billion specifically so it can be used. Read for the word \textbf{earning}.
\end{lrtrapbox}

\subsection{The interest-sensitive gap}

The \textbf{IS gap} is one measure of interest rate risk. It is calculated for assets and liabilities based on \textbf{the time period over which their interest rates adjust to changes in market rates} --- for items that reprice over a specific period, called a \term{bucket}.

\begin{lrfmlbox}[Three forms of the interest-sensitive gap]
\[ \text{Dollar IS gap} \;=\; \text{ISA} - \text{ISL} \]
\[ \text{Relative IS gap} \;=\; \frac{\text{Dollar IS gap}}{\text{Size of the financial institution}} \qquad\qquad
   \text{Interest sensitivity ratio (ISR)} \;=\; \frac{\text{ISA}}{\text{ISL}} \]
\vk{ISA = interest-sensitive (repriceable) assets; ISL = interest-sensitive liabilities. The dollar gap is a difference, the relative gap divides that difference by institution size, and the ISR is a \emph{ratio} of the two levels --- so its neutral value is \textbf{1}, not zero.}
\end{lrfmlbox}

\begin{center}
\small
\begin{tabularx}{\textwidth}{@{}X X@{}}
\toprule
\rowcolor{PaleGreen}\textbf{An asset-sensitive financial firm has:} & \cellcolor{PaleRed}\textbf{A liability-sensitive financial firm has:} \\
\midrule
\textbf{Positive} dollar IS gap \ (ISA $>$ ISL) & \textbf{Negative} dollar IS gap \ (ISA $<$ ISL) \\
\textbf{Positive} relative IS gap & \textbf{Negative} relative IS gap \\
Interest sensitivity ratio \textbf{greater than one} & Interest sensitivity ratio \textbf{less than one} \\
\bottomrule
\end{tabularx}
\end{center}

\begin{center}
\begin{tikzpicture}[
  hd/.style={draw=FRMGold, fill=PaleGold, rounded corners=1pt, align=center,
             font=\scriptsize\sffamily\bfseries, text width=3.4cm, inner sep=4pt},
  gpos/.style={draw=FRMGreen, fill=PaleGreen, rounded corners=1pt, align=center,
             font=\scriptsize\sffamily, text width=4.3cm, inner sep=4pt},
  gneg/.style={draw=FRMRed, fill=PaleRed, rounded corners=1pt, align=center,
             font=\scriptsize\sffamily, text width=4.3cm, inner sep=4pt},
  ar/.style={-{Stealth[length=4.5pt]}, FRMNavy, thick}
]
\node[hd] (top) at (6.6,2.6) {IS gap \\ $=$ ISA $-$ ISL};
\node[gpos] (p1) at (3.0,1.0) {\textbf{Positive IS gap} \ (asset sensitive)\\ ISA $>$ ISL};
\node[gneg] (n1) at (10.2,1.0) {\textbf{Negative IS gap} \ (liability sensitive)\\ ISA $<$ ISL};
\node[gpos] (p2) at (3.0,-0.9) {Rates \textbf{rise} $\rightarrow$ NII and NIM \textbf{up} \\ Rates \textbf{fall} $\rightarrow$ NII and NIM \textbf{down}};
\node[gneg] (n2) at (10.2,-0.9) {Rates \textbf{rise} $\rightarrow$ NII and NIM \textbf{down} \\ Rates \textbf{fall} $\rightarrow$ NII and NIM \textbf{up}};
\node[gpos] (p3) at (3.0,-2.9) {\textbf{Aggressive action:} buy rate-sensitive assets, sell rate-sensitive liabilities};
\node[gneg] (n3) at (10.2,-2.9) {\textbf{Aggressive action:} sell rate-sensitive assets, buy rate-sensitive liabilities};
\draw[ar] (top.south west) -- (p1.north);
\draw[ar] (top.south east) -- (n1.north);
\draw[ar] (p1) -- (p2);
\draw[ar] (n1) -- (n2);
\draw[ar] (p2) -- (p3);
\draw[ar] (n2) -- (n3);
\end{tikzpicture}
\figcap{The single most polarity-dense object in the area: two gap signs, two rate directions, two margin outcomes and two prescribed actions, all mirror images. Hold one anchor --- \textbf{positive gap means assets reprice first, so rising rates help} --- and derive the other seven from it rather than memorising eight facts.}
\end{center}

\subsection{Two approaches to IS gap management}

\begin{center}
\small
\begin{tabularx}{\textwidth}{@{}p{3.3cm} X@{}}
\toprule
\rowcolor{FRMSoft}\textbf{Approach} & \textbf{What it does} \\
\midrule
\textbf{1. Defensive gap management} &
Set the interest-sensitive gap \textbf{as close to zero as possible}, to reduce the expected volatility of net interest income.
\[ \Delta\text{NII} = \begin{pmatrix}\text{Overall change in}\\ \text{interest rate (\% points)}\end{pmatrix} \times \begin{pmatrix}\text{Cumulative gap}\\ \text{(in dollars)}\end{pmatrix} \]
where the cumulative gap is the \textbf{sum of all IS gaps across buckets}. Note: \textbf{a zero gap does not eliminate all interest rate risk}, because the rates attached to assets and liabilities are not perfectly correlated in the real world. \\
\addlinespace
\textbf{2. Aggressive gap management} &
Deliberately position the gap according to the rate forecast. \\
\bottomrule
\end{tabularx}
\end{center}

\begin{center}
\small
\begin{tabularx}{\textwidth}{@{}p{3.6cm} p{3.2cm} X@{}}
\toprule
\rowcolor{FRMSoft}\textbf{Expected change in interest rates (management's forecast)} & \textbf{Best IS gap position to be in} & \textbf{Aggressive management's most likely action} \\
\midrule
\textbf{Rising} market interest rates & \textbf{Positive} IS gap & Increase interest-sensitive assets; decrease interest-sensitive liabilities \\
\addlinespace
\textbf{Falling} market interest rates & \textbf{Negative} IS gap & Decrease interest-sensitive assets; increase interest-sensitive liabilities \\
\bottomrule
\end{tabularx}
\end{center}

\begin{lrkeybox}[Eliminating an interest-sensitive gap]
\begin{center}
\scriptsize
\begin{tabularx}{0.97\textwidth}{@{}p{2.9cm} p{3.4cm} X@{}}
\toprule
\rowcolor{FRMSoft}\textbf{Position} & \textbf{The risk} & \textbf{Possible management responses} \\
\midrule
\textbf{Positive gap} \newline ISA $>$ ISL (asset sensitive) &
Losses if interest rates \textbf{fall}, because the net interest margin will be reduced. &
1. Do nothing (perhaps rates will rise or be stable). \newline
2. \textbf{Extend} asset maturities or \textbf{shorten} liability maturities. \newline
3. Increase interest-sensitive liabilities or reduce interest-sensitive assets. \\
\addlinespace
\textbf{Negative gap} \newline ISA $<$ ISL (liability sensitive) &
Losses if interest rates \textbf{rise}, because the net interest margin will be reduced. &
1. Do nothing (perhaps rates will fall or be stable). \newline
2. \textbf{Shorten} asset maturities or \textbf{lengthen} liability maturities. \newline
3. Decrease interest-sensitive liabilities or increase interest-sensitive assets. \\
\bottomrule
\end{tabularx}
\end{center}
\end{lrkeybox}

\begin{lrtrapbox}[``Do nothing'' is a legitimate offered response]
Both rows list \textbf{do nothing} as response 1, on the grounds that the forecast may not materialise. \S5B.4 notes that GARP's correct answer is sometimes \textbf{no action}, and this table is the reading that licenses it. Do not eliminate an option because it proposes inaction.
\end{lrtrapbox}

\begin{lrtrapbox}[Maturity adjustments run opposite to the gap sign]
To close a \textbf{positive} gap: \textbf{extend} asset maturities (making assets \emph{less} rate-sensitive) or \textbf{shorten} liability maturities. To close a \textbf{negative} gap: \textbf{shorten} assets, \textbf{lengthen} liabilities. Longer maturity means \emph{less} interest-sensitive, and that inverse is the step people skip.
\end{lrtrapbox}

% =====================================================================
\lo{LTR-17 c}{Describe duration gap management and apply this strategy to protect a bank's net worth.}

\begin{lrdefbox}[Duration, in this reading's terms]
A \textbf{time-weighted measure of maturity} that considers the amount and timing of each of an asset's or liability's cash flows. It is also a measure of interest rate sensitivity --- think Macaulay duration.
\end{lrdefbox}

Changing interest rates can do serious damage to another aspect of a firm's performance --- its \textbf{net worth}, the value of the stockholders' investment in the institution. \textbf{For most companies net worth is more important than their net interest margin.} That requires a different tool.

\begin{lrfmlbox}[Net worth and its change]
\[ \text{NW} = A - L \qquad\qquad \Delta\text{NW} = \Delta A - \Delta L \]
\vk{$A$ = market value of assets; $L$ = market value of liabilities. As market rates change, the value of both changes, and net worth moves by the difference.}
\end{lrfmlbox}

\begin{lrkeybox}[What portfolio theory contributes]
\begin{enumerate}
\item A rise in market rates will cause the market value (price) of \textbf{both} fixed-rate assets \textbf{and} liabilities to decline.
\item \textbf{The longer the maturity} of a firm's assets and liabilities, the more they will tend to decline in market value when rates rise.
\end{enumerate}
So a change in net worth due to changing rates varies with the \emph{relative} maturities of assets and liabilities. A firm with \textbf{longer-duration assets than liabilities} will suffer a greater decline in net worth when rates rise.
\end{lrkeybox}

\subsection{The duration gap, and why it must be leverage-adjusted}

\begin{lrfmlbox}[Duration gap and leverage-adjusted duration gap]
\[ \text{Duration gap} \;=\; \underbrace{D_A}_{\substack{\text{dollar-weighted duration}\\ \text{of the asset portfolio}}} - \underbrace{D_L}_{\substack{\text{dollar-weighted duration}\\ \text{of the liability portfolio}}} \]
\[ \text{Leverage-adjusted duration gap} \;=\; D_A - D_L \times \frac{L}{A} \]
\vk{Because the dollar volume of assets usually exceeds the dollar volume of liabilities --- otherwise the firm would be insolvent --- a firm seeking to minimize the effect of rate fluctuations must \textbf{adjust for leverage}. \textbf{The goal of duration gap management is to set the \emph{leverage-adjusted} duration gap to zero.}}
\end{lrfmlbox}

\begin{lrtrapbox}[The unadjusted gap is the intermediate-result trap]
$D_A - D_L$ and $D_A - D_L(L/A)$ are different numbers, and the first one is always offered. Immunization requires the \textbf{leverage-adjusted} gap to be zero, which means $D_A = D_L \times (L/A)$ --- so at immunization the \emph{unadjusted} gap is \textbf{positive}, not zero. An option asserting that a fully hedged bank has $D_A = D_L$ is wrong for every solvent bank.
\end{lrtrapbox}

\begin{center}
\small
\begin{tabularx}{\textwidth}{@{}p{5.6cm} p{2.6cm} X@{}}
\toprule
\rowcolor{FRMSoft}\textbf{If the leverage-adjusted duration gap is:} & \textbf{And if interest rates:} & \textbf{The institution's net worth will:} \\
\midrule
\textbf{Positive} \ ($D_A > D_L \times L/A$) & Rise \newline Fall & \textbf{Decrease} \newline \textbf{Increase} \\
\addlinespace
\textbf{Negative} \ ($D_A < D_L \times L/A$) & Rise \newline Fall & \textbf{Increase} \newline \textbf{Decrease} \\
\addlinespace
\textbf{Zero} \ ($D_A = D_L \times L/A$) & Rise \newline Fall & \textbf{No change} \newline \textbf{No change} \\
\bottomrule
\end{tabularx}
\end{center}
\figcap{The mechanism behind the signs: with a \textbf{positive} gap, a parallel rate change moves \emph{liability} values by \emph{less} than asset values, so a rate rise drops assets further than liabilities and net worth falls. With a \textbf{negative} gap the reverse holds --- liabilities move more, so a rate rise \emph{increases} net worth. Zero gap immunizes.}

\gapbadge
\gapnote{Both layers give the directional table above but neither prints the equation that produces it. The formula and worked example below are written from the GARP curriculum (Rose and Hudgins, Chapter 7). It is on the \S5B.10 formula-recall watchlist.}

\begin{lrgapbox}
The objective says \emph{apply} the strategy to protect net worth. Applying it requires the equation that turns a duration gap and a rate move into a dollar change in net worth.
\end{lrgapbox}

\begin{lrfmlbox}[Change in the market value of net worth]
\[ \Delta\text{NW} \;=\; \left[ -D_A \times \frac{\Delta r}{1+r} \times A \right] \;-\; \left[ -D_L \times \frac{\Delta r}{1+r} \times L \right] \]
\vk{$A$ = total assets; $D_A$ = average duration of assets; $r$ = the initial interest rate; $\Delta r$ = the change in interest rates; $L$ = total liabilities; $D_L$ = average duration of liabilities. The \textbf{negative signs} are the reminder that prices and rates move in opposite directions. The underlying price relation is $\Delta P / P = -D \times \Delta r / (1+r)$.}
\end{lrfmlbox}

\begin{lrexbox}[The dollar cost of a 2-point rate rise]
A bank has total assets of \textbf{\$300 million} with an average asset duration of \textbf{3.047 years}, and total liabilities of \textbf{\$275 million} with an average liability duration of \textbf{2.669 years}. Interest rates on both assets and liabilities rise from \textbf{8\% to 10\%}. By how much does net worth change?
\tcblower
\textbf{Step 1 --- the common rate factor.}
\[ \frac{\Delta r}{1+r} = \frac{0.02}{1.08} = 0.018519 \]
Note the denominator uses the \textbf{initial} rate of 8\%, not the new 10\%.

\textbf{Step 2 --- the asset leg.}
\[ -3.047 \times 0.018519 \times \$300\text{m} = -\$16.93\text{m} \]

\textbf{Step 3 --- the liability leg.}
\[ -2.669 \times 0.018519 \times \$275\text{m} = -\$13.59\text{m} \]

\textbf{Step 4 --- net worth is assets minus liabilities.}
\[ \Delta\text{NW} = -16.93 - (-13.59) = \mathbf{-\$3.34\text{ million}} \]

\textbf{Cross-check against the gap.} The leverage-adjusted duration gap is
\[ 3.047 - 2.669 \times \frac{275}{300} = 3.047 - 2.447 = +0.600 \]
Positive gap, rates rise, net worth falls --- consistent with the table above. \textbf{Use the sign of the gap as the execution check on the sign of the answer.}

\textbf{What this means.} Both sides of the balance sheet lost value; the bank is worse off only because \textbf{assets lost more}. \$16.93m and \$13.59m are both large and nearly cancel --- the entire exposure is the \$3.34m residual, which is why small duration mismatches on large balance sheets matter.
\end{lrexbox}

\begin{lrtrapbox}[Four ways this calculation is corrupted]
\begin{itemize}
\item \textbf{The wrong $r$ in the denominator.} Use the \textbf{initial} 8\%. Using 10\% gives $-\$3.28$m.
\item \textbf{Dropping the second negative sign.} $-16.93 + (-13.59) = -\$30.52$m. The liability leg is \emph{subtracted}, and it is itself negative, so it comes back as a positive offset.
\item \textbf{Intermediate results.} \$16.93m and \$13.59m are both legitimate quantities in the chain and both will be offered against the correct \$3.34m.
\item \textbf{Unadjusted gap.} $3.047 - 2.669 = 0.378$ versus the leverage-adjusted $0.600$. Both are plausible-looking gap figures and only the second governs immunization.
\end{itemize}
\end{lrtrapbox}

\subsection{Two approaches to duration gap management}

\begin{center}
\small
\begin{tabularx}{\textwidth}{@{}p{3.2cm} p{4.3cm} X@{}}
\toprule
\rowcolor{FRMSoft}\textbf{Expected change in interest rates} & \textbf{Management action} & \textbf{Possible outcome} \\
\midrule
\textbf{Rates will rise} & Reduce $D_A$ and increase $D_L$ \newline (moving closer to a \textbf{negative} duration gap) & Net worth \textbf{increases}, \emph{if} management's rate forecast is correct. \\
\addlinespace
\textbf{Rates will fall} & Increase $D_A$ and reduce $D_L$ \newline (moving closer to a \textbf{positive} duration gap) & Net worth \textbf{increases}, \emph{if} management's rate forecast is correct. \\
\bottomrule
\end{tabularx}
\end{center}
\figcap{This is the \textbf{aggressive} duration gap strategy; the \textbf{defensive} strategy simply targets a leverage-adjusted gap of zero and takes no view. Note that both rows promise the same outcome --- net worth increases --- conditional on the forecast being right. The conditional is the examinable part: aggressive gap management converts interest rate risk into \emph{forecast} risk, it does not remove risk.}

\begin{lrtrapbox}[Rising rates call for a negative gap on both tools, but for different reasons]
Under \textbf{IS gap} management, a rising-rate forecast calls for a \textbf{positive} gap (assets reprice up faster, NIM rises). Under \textbf{duration gap} management, a rising-rate forecast calls for a \textbf{negative} gap (liabilities fall in value faster, net worth rises). \textbf{The two prescriptions have opposite signs for the same forecast}, because one protects a flow and the other protects a value. This is the single highest-value confusable pair in the reading --- add it to the registry.
\end{lrtrapbox}

% =====================================================================
\lo{LTR-17 d}{Discuss the limitations of interest-sensitive gap management and duration gap management.}

\begin{center}
\small
\begin{tabularx}{\textwidth}{@{}p{0.5cm} X@{}}
\toprule
\rowcolor{PaleNavy}\multicolumn{2}{@{}l}{\textbf{Limitations of interest-sensitive gap management}} \\
\midrule
a) & Rates on assets and liabilities \textbf{may not move together}. \\
b) & Rates may change at \textbf{different speeds within a maturity bucket}. \\
c) & \textbf{The point at which assets and liabilities reprice is not easy to identify} --- banks may adjust rates on interest-bearing demand deposits and savings accounts at any time. \\
d) & An aggressive approach to NII and NIM management via IS gap may \textbf{increase earnings volatility}, which may in turn negatively affect the bank's stock price. \\
\bottomrule
\end{tabularx}
\end{center}

\begin{center}
\small
\begin{tabularx}{\textwidth}{@{}p{0.5cm} X@{}}
\toprule
\rowcolor{PaleGold}\multicolumn{2}{@{}l}{\textbf{Limitations of duration gap management}} \\
\midrule
a) & It may \textbf{not be possible to match the durations} of assets to liabilities, even if managers prefer a zero duration gap. Finding assets and liabilities of the same duration that fit the portfolio is often a frustrating task. \\
b) & The durations of \textbf{shorter-term} securities and loans are \textbf{closer to their maturities} than the durations of longer-term securities and loans. \\
c) & It is often difficult to \textbf{define the pattern of cash flows} for instruments such as demand deposits and savings accounts, making duration hard to calculate. \\
d) & Duration gap is more effective for \textbf{small} rate changes and for changes that affect the \textbf{entire yield curve in parallel}, than for large changes or slope changes. \textbf{Convexity} must be added to the calculation for more precise results on large movements. \\
e) & \textbf{Duration itself shifts as rates change.} \\
\bottomrule
\end{tabularx}
\end{center}

\begin{lrtrapbox}[One limitation is shared, and it is the diagnostic one]
Both tools founder on the same instruments: \textbf{demand deposits and savings accounts}. For IS gap the problem is that \emph{repricing dates cannot be identified}; for duration gap it is that \emph{cash flow patterns cannot be defined}. Same product, two different failures --- a distractor attaching the wrong failure to the right tool is the precise build this pairing invites.
\end{lrtrapbox}

\begin{lrtrapbox}[Consolidated trap summary — LTR-17]
\begin{itemize}
\item \textbf{Role.} IS gap protects the \textbf{net interest margin} (a flow). Duration gap protects \textbf{net worth} (a value). Identify the branch before answering.
\item \textbf{Input twin.} NIM divides by \textbf{earning} assets, not total assets. In the worked example, \$45bn not \$50bn.
\item \textbf{Polarity.} \textbf{Positive} IS gap $=$ asset sensitive $=$ ISR $>$ 1; rates \emph{rise} $\rightarrow$ NII and NIM \emph{up}. Negative gap mirrors it.
\item \textbf{Polarity.} \textbf{Positive} leverage-adjusted duration gap; rates \emph{rise} $\rightarrow$ net worth \emph{falls}. Negative gap mirrors it. Zero gap immunizes.
\item \textbf{Sibling (highest value in the reading).} A rising-rate forecast calls for a \textbf{positive IS gap} but a \textbf{negative duration gap}. Opposite signs, same forecast.
\item \textbf{Calculation.} Leverage-adjusted duration gap $= D_A - D_L \times (L/A)$. Immunization means $D_A = D_L(L/A)$, so the \emph{unadjusted} gap is positive at immunization --- $D_A = D_L$ is wrong for any solvent bank.
\item \textbf{Calculation.} $\Delta \text{NW}$ uses $\Delta r / (1 + r)$ with the \textbf{initial} $r$, and the liability leg is \textbf{subtracted}.
\item \textbf{Intermediate result.} \$16.93m (asset leg), \$13.59m (liability leg) and 0.378 (unadjusted gap) are all offered against \$3.34m and $+0.600$.
\item \textbf{Polarity.} To close a positive IS gap: \textbf{extend} assets or \textbf{shorten} liabilities. Longer maturity $=$ \emph{less} rate-sensitive.
\item \textbf{Scope.} A \textbf{zero IS gap does not eliminate all interest rate risk}, because asset and liability rates are not perfectly correlated.
\item \textbf{Scope.} ``\textbf{Do nothing}'' is one of the three listed responses to a non-zero gap. Do not eliminate an inaction option on sight.
\item \textbf{Sibling.} Demand deposits and savings accounts defeat \textbf{IS gap} via unidentifiable \emph{repricing dates}, and \textbf{duration gap} via undefinable \emph{cash flow patterns}.
\end{itemize}
\end{lrtrapbox}
